% sections/shpn_formalism.tex
% Final section: introduces Signal Hierarchical Petri Nets as a named formalism.
% Absorbs background + model construction from the old draft.
% No internal phase labels, model filenames, or run IDs.

\subsection{Petri Nets as a Modelling Framework for Molecular Networks}

Petri nets provide a natural formalism for modelling concurrent molecular
processes. In the biological Petri net framework~\cite{heiner2008}, chemical
species are represented as places holding real-valued tokens (concentrations),
and chemical reactions are transitions that consume and produce tokens according
to stoichiometry. Test arcs allow transitions to read a concentration without
consuming it, modelling catalytic effects such as enzyme action. Standard Petri
net semantics accommodate deterministic (continuous), stochastic (discrete), and
hybrid interpretations within the same structural specification, making the
formalism well-suited for models that must span multiple scales of organisation
and copy number.

\subsection{Signal Hierarchical Petri Nets}
\label{sec:shpn_def}

Signal Hierarchical Petri Nets (SHPN)~\cite{shypn2026} extend the standard
biological Petri net with three additional constructs.

\emph{Signal places} are information channels carrying regulatory content rather
than material substance. Unlike ordinary places, signal places are not depleted
by arc consumption: a signal can be read by multiple downstream transitions
simultaneously, representing broadcast regulatory effects such as a transcription
factor activating many target genes at once.

\emph{Signal flow arcs} consume tokens from signal places during transition
firing, providing a formal representation of signal consumption in commitment
circuits. A signal flow arc from a signal place $s$ to transition $t$ means that
$t$ can fire only while $s$ carries a marking above a threshold, and that firing
reduces $s$ by the arc weight. This mechanism models the coupling between
upstream signal state and downstream reaction availability that underlies
commitment cascade preemption.

\emph{Adaptive transitions} parameterise their firing rates as explicit functions
of current place markings through user-defined kinetic expressions. Standard
mass-action, Michaelis--Menten saturation, Hill-type cooperative terms, and
pH-dependent inhibition constants are all expressible within the adaptive
transition rate language. This kinetic flexibility is essential for representing
saturable receptor binding, cooperative transcription factor interactions,
energy-coupled phosphorylation, and the pH-Hill cross-inhibition mechanism
(equation~\ref{eq:pHhill}) used in this work.

The stochastic semantics of Signal Hierarchical Petri Nets follow a
continuous-time Markov chain identical in structure to the Gillespie stochastic
simulation algorithm~\cite{gillespie1977}: each transition fires at a rate
proportional to its kinetic expression evaluated at the current marking, with
the next firing event drawn from an exponential distribution and the firing
transition selected proportionally to its rate weight. Molecular noise therefore
emerges naturally from the simulation dynamics without requiring explicit noise
injection. The framework is intrinsically hybrid: when copy numbers are large, the
Markov chain approximates deterministic kinetics; when copy numbers are small, it
enters the high-noise regime in which individual stochastic transitions produce
proportionally large perturbations.

\subsection{Model Architecture}
\label{sec:model_arch}

The GATA1/PU.1 commitment model consists of 28 places and 29 transitions,
partitioned into six biological compartments: extracellular cytokines (EPO,
GCSF); membrane receptors (EPO receptor and GCSF receptor, together with their
bound, internalised, and recycling-pool forms); cytoplasmic and nuclear mRNA for
both GATA1 and PU.1; cytoplasmic and nuclear protein for both factors;
phosphorylated nuclear GATA1; and the energy currency species (ATP, ADP, GTP,
GDP, inorganic phosphate). The model topology is illustrated in
Fig.~\ref{fig:model}.

\begin{figure}[!ht]
  \centering
  \includegraphics[width=\linewidth]{fig_model_topology_v6}
  \caption{\textbf{Signal Hierarchical Petri Net topology of the GATA1/PU.1
  commitment model.} 28 places and 29 transitions spanning six biological
  compartments. Filled circles are ordinary places; double circles are signal
  places; boxes are adaptive transitions. Arc annotations show stoichiometric
  weights; signal flow arcs are shown as dashed lines.}
  \label{fig:model}
\end{figure}

Transitions represent cytokine receptor binding and unbinding, receptor
internalisation and recycling, mRNA transcription, nuclear export, mRNA
degradation, protein translation, nuclear import, protein degradation, GATA1
phosphorylation, and ubiquitin-mediated PU.1 degradation. Energy metabolism is
represented through ATP synthesis from ADP and inorganic phosphate, and GTP
regeneration from GDP and inorganic phosphate, so that the energy species are
dynamically coupled to the signalling network rather than held at fixed
concentrations.

The model implements two opposing regulatory arms. In the erythroid arm, bound
EPO receptor promotes GATA1 transcription; the resulting GATA1 protein is
exported, translated, imported into the nucleus, and phosphorylated by a kinase
whose activity scales with EPO receptor occupancy. Phosphorylated nuclear GATA1
then drives PU.1 ubiquitin-mediated degradation through a transition whose rate
is proportional to the product of phosphorylated GATA1 abundance and nuclear PU.1
abundance, representing GATA1-dependent recruitment of an E3 ubiquitin ligase to
PU.1. In the myeloid arm, bound GCSF receptor promotes PU.1 transcription;
nuclear PU.1 suppresses GATA1 transcription through a competitive inhibition term.
These kinetic forms produce the double-negative feedback topology necessary and
sufficient for bistability~\cite{huang2007}.

Consistent with the established biology that uncommitted progenitors carry a
slight GATA1 bias~\cite{cantor2002}, the GATA1 basal transcription rate is set
33\% above the PU.1 rate (0.08 versus 0.06~mM/s). This asymmetry is a structural
modelling assumption reflecting the observed GATA1-primed state of common myeloid
progenitors; it is not a calibration artefact.

\subsection{Kinetic Parameters and Calibration}
\label{sec:calibration}

Rate constants were derived from published measurements for EPO receptor binding
kinetics~\cite{huang2007}, GATA1 and PU.1 transcription and degradation
rates~\cite{chickarmane2009}, and receptor internalisation and recycling dynamics.
Energy metabolism parameters were calibrated to maintain physiologically
reasonable ATP and GTP steady states
($\text{ATP}+\text{ADP} \approx 33\,\text{mM}$;
$\text{GTP}+\text{GDP} \approx 17.73\,\text{mM}$), consistent with published measurements for
erythroid progenitor cells. Cross-repression strengths and basal transcription
rates were calibrated to produce, at population-scale conditions, a sharp
commitment threshold near the EPO* values observed in bulk colony assays.

\subsection{Single-Cell Copy Number Scaling}
\label{sec:scaling}

To access the stochastic regime corresponding to physiological single-cell copy
numbers, the initial markings of all intracellular molecular species (mRNA,
cytoplasmic and nuclear proteins, energy metabolites) were set to one-hundredth
of their population-scale reference values. Extracellular cytokines and membrane
receptor species were not scaled, preserving concentrations set by the bone
marrow niche and actual receptor numbers per cell. At one-hundredth of standard
conditions, the nuclear GATA1 place has an initial marking of approximately eight
molecules, consistent with published estimates of master transcription factor
copy numbers in haematopoietic progenitor cells~\cite{chang2008}.

At these copy numbers, the Gillespie stochastic semantics of the framework
simulate molecular events as discrete random occurrences. A single transcription
event adds one molecule to a place containing eight, producing a relative
fluctuation of twelve percent. This is the mechanistic origin of the high
cell-to-cell fate variability described in the Results: the noise is not injected
artificially but emerges from the discreteness of molecular events in the
decision layer.

Bimolecular transitions involving one intracellular species and one receptor
species were treated with care: transition rates were expressed in terms of the
scaled intracellular species abundance while receptor occupancy was computed from
the unscaled receptor concentration, ensuring that receptor-coupled signalling
was physiologically consistent rather than artificially attenuated.

\subsection{Comparison with Deterministic Bistable Models}

The GATA1/PU.1 mutual antagonism circuit belongs to the class of toggle-switch
networks whose deterministic dynamics produce bistability through mutual repression
and self-activation~\cite{huang2007}. In the deterministic limit---large copy
numbers, negligible fluctuations---the present model recovers the same two-attractor
structure described by Huang \etal~\cite{huang2007} and Chickarmane
\etal~\cite{chickarmane2009}. The advance provided by the Signal Hierarchical
Petri Net implementation is the ability to evaluate the same circuit across the
full copy-number spectrum, from the molecular-count noise of a single progenitor
to the near-deterministic dynamics of a bulk colony, within a single model
specification. The layer decomposition revealed in the Results---and in particular
the identification of the decision layer as the exclusive locus of stochastic
fate specification---is an emergent finding of the stochastic analysis that is
invisible to deterministic models. Signal Hierarchical Petri Nets provide the
formalism that makes this analysis possible.
